% Generated deterministically from the audited Markdown source.
% Responsible author: Carptopus.
\documentclass[11pt]{article}
\usepackage[a4paper,margin=27mm]{geometry}
\usepackage[T1]{fontenc}
\usepackage[utf8]{inputenc}
\usepackage{lmodern}
\usepackage{microtype}
\usepackage{amsmath,amssymb,amsthm,mathtools}
\usepackage{booktabs,tabularx,array,graphicx}
\usepackage{needspace}
\usepackage{xcolor}
\usepackage{fancyvrb}
\usepackage[colorlinks=true,linkcolor=blue!55!black,citecolor=blue!55!black,urlcolor=blue!65!black]{hyperref}
\usepackage{fancyhdr}

\newtheorem{theorem}{Theorem}[section]
\newtheorem{proposition}[theorem]{Proposition}
\newtheorem{lemma}[theorem]{Lemma}
\numberwithin{equation}{section}

\pagestyle{fancy}
\fancyhf{}
\fancyhead[L]{Carptopus}
\fancyhead[R]{Two cyclotomic minimum-degree cases}
\fancyfoot[C]{\thepage}
\setlength{\headheight}{14pt}
\setlength{\parindent}{0pt}
\setlength{\parskip}{0.45em}
\setlength{\emergencystretch}{4em}
\urlstyle{same}

\title{Exact certificates for the first two cases not covered\\
by Steinberger's cyclotomic minimum-degree theorem}
\author{Carptopus\\\href{mailto:carptopus@163.com}{carptopus@163.com}}
\date{27 August 2026}

\hypersetup{
  pdftitle={Exact certificates for the first two cases not covered by Steinberger's cyclotomic minimum-degree theorem},
  pdfauthor={Carptopus},
  pdfsubject={Exact integer certificates for the cyclotomic nonnegative-coefficient minimum-degree problem at n=46189 and n=96577},
  pdfkeywords={cyclotomic polynomial, nonnegative coefficients, minimum degree, exact certificate, CRT array, linear programming}
}

\begin{document}
\maketitle

\begin{center}
\fcolorbox{black!35}{black!4}{\parbox{0.88\textwidth}{\centering\small
Internal candidate manuscript, version 0.1-beta. The mathematical proof and
complete Markdown manuscript have passed project-internal zero-trust audits.
External mathematical review and formal peer review are pending.}}
\end{center}

\begin{abstract}


Let $\Phi_n(x)$ be the $n$-th cyclotomic polynomial. Steinberger asked for the
lowest degree of a nonzero polynomial with nonnegative coefficients divisible
by $\Phi_n$ and conjectured that the unique monic minimizer is the regular
$p$-gon polynomial

\nopagebreak[4]
\[
1+x^{n/p}+\cdots+x^{(p-1)n/p},
\]

where $p$ is the smallest prime divisor of $n$. His general theorem leaves
$46189=11\cdot13\cdot17\cdot19$ and
$96577=13\cdot17\cdot19\cdot23$ as the first two uncovered values. We prove
the conjecture for both values. The proof is computer assisted but exact: for
each parameter we provide an integer element of the orthogonal complement of
the cyclotomic-array space which vanishes on the regular polygon and has a
strict sign at every other position up to the proposed boundary degree. The
certificates are checked by two exact routes, one using CRT tensor marginals
and one expanding the full multidimensional array and verifying every fiber
sum. Consequently the minimum degrees are $41990$ and $89148$, respectively,
and the monic minimizers over the nonnegative reals are unique. These are two
fixed-parameter results, not a proof of the general conjecture.

\end{abstract}

\section{Introduction}


For a positive integer $n$, consider the optimization problem

\nopagebreak[4]
\[
\min\{\deg f:0\ne f\in\mathbb R_{\geq0}[x],\ \Phi_n(x)\mid f(x)\}.
\tag{1.1}
\]

If $p$ is the smallest prime divisor of $n$, the polynomial

\nopagebreak[4]
\[
R_{n,p}(x)=1+x^{n/p}+\cdots+x^{(p-1)n/p}
\tag{1.2}
\]

is divisible by $\Phi_n(x)$ and has degree $(p-1)n/p$. Steinberger \cite{steinberger2012}
conjectured that this is always the minimum degree and that $R_{n,p}$ is the
unique monic minimizer. His Theorem 1 proves the conjecture when $n$ is even,
when $n$ is a prime power, and, for squarefree notation
$n=pq_1\cdots q_k$, whenever

\nopagebreak[4]
\[
\frac{2}{p}>\frac1{q_1}+\cdots+\frac1{q_k}.
\tag{1.3}
\]

The same paper identifies

\nopagebreak[4]
\[
46189=11\cdot13\cdot17\cdot19,
\qquad
96577=13\cdot17\cdot19\cdot23
\tag{1.4}
\]

as the first two values not covered by that theorem. It also explains that
each fixed case is a finite linear-programming problem, while noting that the
resulting systems were too large for the computations carried out there.

The purpose of this note is to settle exactly these two fixed cases. Our main
result is the following.

\Needspace{0.28\textheight}
\begin{theorem}

Let $f\in\mathbb R_{\geq0}[x]$ be nonzero.

\begin{enumerate}
\item If $\Phi_{46189}(x)\mid f(x)$, then $\deg f\geq41990$. Equality holds
   precisely for positive scalar multiples of
   $1+x^{4199}+\cdots+x^{10\cdot4199}$.
\item If $\Phi_{96577}(x)\mid f(x)$, then $\deg f\geq89148$. Equality holds
   precisely for positive scalar multiples of
   $1+x^{7429}+\cdots+x^{12\cdot7429}$.

\end{enumerate}

In particular, in each case the displayed polynomial is the unique monic
minimizer.

\end{theorem}

The mathematical reduction is Steinberger's separation criterion \cite[Lemma~1]{steinberger2012}. The new finite objects are two exact integer certificates. Floating
point linear programming was used to discover them, but it is not part of the
proof: the final claims depend only on the integer files and exact arithmetic.

Section 2 recalls the separation criterion and makes the coefficient field in
the uniqueness statement explicit. Section 3 gives the CRT-array
representation and the certificate data. Section 4 describes the exact
verification. Section 5 records how the larger certificate was found, while
separating that history from the proof. Section 6 states the scope relative to
prior work.

\section{The separation criterion}


We use Steinberger's notation in a form adapted to exact certificates. Let $n$
be squarefree, let

\nopagebreak[4]
\[
n=p_1p_2\cdots p_s,
\qquad p_1=p<p_2<\cdots<p_s,
\]

and put

\nopagebreak[4]
\[
P=\frac np,
\qquad d=(p-1)P.
\]

Let $V_n\subseteq\mathbb R^n$ be the space of coefficient vectors of
polynomials of degree less than $n$ which are divisible by $\Phi_n$. Define

\nopagebreak[4]
\[
H=\{0,P,2P,\ldots,(p-1)P\}
\tag{2.1}
\]

and

\nopagebreak[4]
\[
F^+=\{0,1,\ldots,d-1\}\setminus H.
\tag{2.2}
\]

Thus $H$ is the support of the regular $p$-gon polynomial, including its
degree-$d$ endpoint. Steinberger's Lemma 1 uses a vector in $V_n^\perp$ which
is zero on $H$ and strictly positive on $F^+$. We use the negative of that
vector.

\Needspace{0.28\textheight}
\begin{proposition}[Steinberger's criterion, sign reversed]

Suppose there is a vector $w\in V_n^\perp$ such that

\nopagebreak[4]
\[
w_j=0\quad(j\in H),
\qquad
w_j<0\quad(j\in F^+).
\tag{2.3}
\]

Then every nonzero $f\in\mathbb R_{\geq0}[x]$ divisible by $\Phi_n$ and of
degree at most $d$ is supported on $H$.

\end{proposition}

\begin{proof}

Pad the coefficient vector $a$ of $f$ with zeros to length $n$. Because
$\deg f\leq d<n$ and $\Phi_n\mid f$, we have $a\in V_n$. Hence

\nopagebreak[4]
\[
0=\langle w,a\rangle=\sum_{j=0}^{d}w_ja_j.
\]

Every summand is nonpositive, and every summand indexed by $F^+$ is strictly
negative whenever the corresponding coefficient is positive. Therefore
$a_j=0$ for every $j\in F^+$, so the support is contained in $H$.
\end{proof}


For completeness, we record the step which turns support containment into the
real-coefficient uniqueness statement. It is important here not to appeal
merely to irreducibility over $\mathbb Q$, since our coefficients lie in
$\mathbb R$.

\Needspace{0.28\textheight}
\begin{lemma}[real-coefficient boundary rigidity]

Let $n$ be squarefree, let $p$ be its smallest prime divisor, and put $P=n/p$.
Suppose $0\ne f\in\mathbb R_{\geq0}[x]$, $\Phi_n\mid f$,
$\deg f\leq(p-1)P$, and $\operatorname{supp}(f)\subseteq H$. Then

\nopagebreak[4]
\[
f(x)=c\bigl(1+x^P+\cdots+x^{(p-1)P}\bigr)
\tag{2.4}
\]

for some $c>0$.

\end{lemma}

\begin{proof}

Write $f(x)=g(x^P)$ with $g\in\mathbb R[x]$ and $\deg g\leq p-1$. Let
$\zeta_n$ be a primitive $n$-th root of unity. For every
$a\in(\mathbb Z/n\mathbb Z)^*$,

\nopagebreak[4]
\[
0=f(\zeta_n^a)=g\bigl((\zeta_n^P)^a\bigr).
\]

The root $\zeta_n^P$ is primitive of order $p$. Because $n$ is squarefree,
the reduction map

\nopagebreak[4]
\[
(\mathbb Z/n\mathbb Z)^*\longrightarrow(\mathbb Z/p\mathbb Z)^*
\tag{2.5}
\]

is surjective by the Chinese remainder theorem. Consequently $g$ vanishes at
every primitive $p$-th root. Therefore $\Phi_p\mid g$ in $\mathbb R[x]$.
Since both $\deg g\leq p-1$ and $\deg\Phi_p=p-1$, we have
$g=c\Phi_p$ for a real scalar $c$. Nonzero nonnegative coefficients force
$c>0$, proving (2.4).
\end{proof}


Proposition 2.1 and Lemma 2.2 show that a single boundary certificate proves
both the degree lower bound and boundary uniqueness. Indeed, a polynomial of
degree less than $d$ would be forced to equal a positive scalar multiple of
the regular polygon, which has degree exactly $d$.

\section{CRT arrays and the two certificates}


We recall the exact coordinate model used to store the certificates. For
$n=p_1\cdots p_s$, define

\nopagebreak[4]
\[
\psi(t_1,\ldots,t_s)
=\sum_{r=1}^s t_r\frac{n}{p_r}\pmod n,
\qquad
0\leq t_r<p_r.
\tag{3.1}
\]

Every factor $n/p_r$ is invertible modulo $p_r$, so $\psi$ is a bijection. Its
inverse sends an exponent $j$ to

\nopagebreak[4]
\[
t_r(j)=j\left(\frac{n}{p_r}\right)^{-1}\pmod {p_r}.
\tag{3.2}
\]

Under this identification, the space $V_n$ is spanned by indicators of fibers
parallel to the coordinate axes \cite{steinberger2012}. It follows that

\nopagebreak[4]
\[
V_n^\perp
=\mathbf1_{p_1}^{\perp}\otimes\cdots\otimes
  \mathbf1_{p_s}^{\perp}.
\tag{3.3}
\]

For $1\leq i<p_r$, let $b_i^{(r)}=e_i-e_0$. These vectors form an integer
basis of $\mathbf1_{p_r}^{\perp}$. An integer coordinate array

\nopagebreak[4]
\[
U=(U_{i_1,\ldots,i_s}),
\qquad 1\leq i_r<p_r,
\]

therefore determines

\nopagebreak[4]
\[
W=
\sum_{i_1=1}^{p_1-1}\cdots\sum_{i_s=1}^{p_s-1}
U_{i_1,\ldots,i_s}
b_{i_1}^{(1)}\otimes\cdots\otimes b_{i_s}^{(s)}
\in V_n^\perp.
\tag{3.4}
\]

The certificate value at exponent $j$ is

\nopagebreak[4]
\[
w_j=W_{t_1(j),\ldots,t_s(j)}.
\tag{3.5}
\]

The two retained boundary certificates have the following exact summaries.

\begin{center}
\small
\resizebox{\textwidth}{!}{%
\begin{tabular}{r l r l r r}
\toprule
$n$ & prime factors & $P=n/p$ & logical shape of $U$ & $\min_{F^+}w_j$ & $\max_{F^+}w_j$ \\
\midrule
46189 & $11,13,17,19$ & 4199 & $10\times12\times16\times18$ & $-4272093684$ & $-999961$ \\
96577 & $13,17,19,23$ & 7429 & $12\times16\times18\times22$ & $-40866441870$ & $-999967$ \\
\bottomrule
\end{tabular}%
}
\end{center}


For $n=46189$, all 11 values on

\nopagebreak[4]
\[
H=\{0,4199,\ldots,10\cdot4199\}
\]

are zero, and all other 41,980 values through degree 41,990 are negative. For
$n=96577$, all 13 values on

\nopagebreak[4]
\[
H=\{0,7429,\ldots,12\cdot7429\}
\]

are zero, and all other 89,136 values through degree 89,148 are negative.

The certificate files are anchored by SHA-256:

\begin{center}
\small
\resizebox{\textwidth}{!}{%
\begin{tabular}{r l l}
\toprule
$n$ & file & SHA-256 \\
\midrule
46189 & {\small\path{boundary-uniqueness-46189.npz}} & {\small\path{79d9042e4c050a7a2f5f12b95f47b2c53857b55a531419ece063fb6d6dc051b4}} \\
96577 & {\small\path{boundary-uniqueness-96577.npz}} & {\small\path{78a03fa247e2d33cd232e7d0365c767c106feb8895ea55e16ababb980ce4a563}} \\
\bottomrule
\end{tabular}%
}
\end{center}


The NPZ files store the integer coordinates as flat vectors; the verification
programs reshape them to the logical tensor dimensions displayed above.

\begin{proof}[Proof of Theorem 1.1]

For each of the two values of $n$, exact verification establishes
$W\in V_n^\perp$ for the reconstructed array, together with zero values on
$H$ and strict negativity on $F^+$. Proposition 2.1 forces every nonnegative multiple of degree at most the
proposed boundary to be supported on $H$. Lemma 2.2 then shows that it is a
positive scalar multiple of the corresponding regular polygon. This proves
the lower bound, characterizes equality, and gives monic uniqueness.
\end{proof}


\section{Exact verification}


The proof objects are the final integer arrays, not the linear-programming
transcripts. We provide two exact evaluation routes.

\subsection{Tensor-marginal evaluation}


Formula (3.4) can be evaluated without materializing all $n$ entries. For
$t=(t_1,\ldots,t_s)$, let $Z(t)=\{r:t_r=0\}$. A nonzero coordinate $t_r$
forces $i_r=t_r$, while a zero coordinate contributes a factor $-1$ and a sum
over $i_r$. Thus

\nopagebreak[4]
\[
W_t=(-1)^{|Z(t)|}
\sum_{\substack{1\leq i_r<p_r\ (r\in Z(t))\\
i_r=t_r\ (r\notin Z(t))}}
U_{i_1,\ldots,i_s}.
\tag{4.1}
\]

The marginal verifiers use (3.2) and (4.1) to recompute every exponent through
the boundary using exact integer arithmetic. They check $n$, the prime tuple,
the boundary, the support $H$, and the coordinate count, and recompute the
certificate hash for comparison with the accompanying metadata. The fixed
hashes in Section 3 are the manuscript-level anchors.

\subsection{Full-array verification}


The second route explicitly expands (3.4) to arrays of shapes

\nopagebreak[4]
\[
11\times13\times17\times19
\quad\text{and}\quad
13\times17\times19\times23.
\]

It then checks that every fiber sum in each of the four coordinate directions
is exactly zero. This verifies membership in $V_n^\perp$ without calling the
marginal evaluation routine. Finally it applies the CRT map to all boundary
positions and checks the required zero set and strict sign.

All sums are exact integers. A deliberately conservative overflow audit bounds
the largest possible fiber accumulation for the three retained certificates
by $4.95\times10^{16}$, below the signed 64-bit limit
$9.22\times10^{18}$. An additional audit expanded the arrays with arbitrary
precision Python integers and reproduced the same extrema, counts, zero sets,
and fiber sums.

For $n=46189$ we also retain a logically redundant strict-separation
certificate for all degrees at most 41,989. Its SHA-256 is

{\small\path{9f0e34517f3ae92161e359e0345563d1ae4ff73a5cdc99693c0e82657dce0e80}}.

It provides a separate check of the minimum-degree conclusion, although the
boundary certificate alone already proves the full statement.

\section{Certificate discovery and integerization}


This section records how the proof objects were found. None of the numerical
statuses described here is used as a theorem premise.

An element of $V_n^\perp$ was parametrized by its integer-basis coordinates
$U$. The boundary problem imposes $p$ homogeneous equalities on $H$ and strict
inequalities on every position in $F^+$. In floating-point discovery these
were normalized to inequalities of the form $w_j\leq-1$.

For $n=46189$, a direct sparse formulation produced a feasible point which
could be scaled, rounded, and repaired exactly. For $n=96577$, the recorded
discovery run of the original 76,032-variable formulation reached its time
limit without producing a proof object. This was a solver outcome, not a
mathematical conclusion.

The successful reformulation eliminated 12 independent boundary equalities
using a pivot minor of determinant $+1$. The recorded reduced bounded
interior-point run then produced a strict floating-point separator.
After scaling by $10^6$ and rounding, the boundary equalities were repaired
over the integers using a 12-by-12 minor of determinant $-1$. The resulting
integer vector was then discarded as a numerical candidate unless and until
the exact verifiers accepted it.

Metadata fields recording the floating-point scale or repair size describe
this discovery history only. They are not proof evidence and need not be
reconstructed from the final NPZ files.

\section{Scope and relation to prior work}


Steinberger \cite{steinberger2012} formulated the optimization problem, proposed the general
conjecture, established the separation criterion used here, and proved the
large parameter range summarized by (1.3). The same paper explicitly lists
$46189$ and $96577$ as the first two values outside that theorem. The present
note contributes exact certificates and their verification for those two
fixed values.

The result does not prove the conjecture for all squarefree integers with four
prime factors, does not provide an infinite parameter family, and does not
settle the next listed value $215441$. The exact computation also does not by
itself explain why a separator should exist in general.

A targeted search completed on 27 August 2026 found no public paper, preprint,
author update, or indexed computational archive directly treating either of
these two certificates. This supports writing the fixed-parameter result but
is not an assertion of exhaustive global priority. Unindexed manuscripts and
independent unpublished computations remain possible.

\section{Reproducibility}


The exact certificate files, metadata, and verification programs are retained
in the project {\small\path{verification}} directory. With Python, NumPy, and these files
available, the principal checks are:

\begin{Verbatim}[fontsize=\footnotesize]
& '.\.venv\Scripts\python.exe' `
  'loops\CYCLOTOMIC-MINDEG-0001\verification\verify_boundary_uniqueness.py' `
  --certificate 'loops\CYCLOTOMIC-MINDEG-0001\verification\results\boundary-uniqueness-46189.npz' `
  --metadata 'loops\CYCLOTOMIC-MINDEG-0001\verification\results\boundary-uniqueness-46189.json'

& '.\.venv\Scripts\python.exe' `
  'loops\CYCLOTOMIC-MINDEG-0001\verification\verify_full_array.py' `
  --mode boundary `
  --certificate 'loops\CYCLOTOMIC-MINDEG-0001\verification\results\boundary-uniqueness-46189.npz' `
  --metadata 'loops\CYCLOTOMIC-MINDEG-0001\verification\results\boundary-uniqueness-46189.json'

& '.\.venv\Scripts\python.exe' `
  'loops\CYCLOTOMIC-MINDEG-0001\verification\verify_boundary_96577.py' `
  --certificate 'loops\CYCLOTOMIC-MINDEG-0001\verification\results\boundary-uniqueness-96577.npz' `
  --metadata 'loops\CYCLOTOMIC-MINDEG-0001\verification\results\boundary-uniqueness-96577.json'

& '.\.venv\Scripts\python.exe' `
  'loops\CYCLOTOMIC-MINDEG-0001\verification\verify_full_array_96577.py' `
  --certificate 'loops\CYCLOTOMIC-MINDEG-0001\verification\results\boundary-uniqueness-96577.npz' `
  --metadata 'loops\CYCLOTOMIC-MINDEG-0001\verification\results\boundary-uniqueness-96577.json'
\end{Verbatim}


The current project environment used Python 3.13.11 and NumPy 2.5.2. SciPy
1.17.0 and HiGHS were used during discovery, but are not required by the
mathematical verification of the final integer arrays.

\section*{Acknowledgement of AI assistance}


The author used OpenAI Codex as an AI-assisted tool for literature discovery,
computational exploration, certificate construction, proof stress testing,
and language editing. The author reviewed the mathematical argument and
assumes responsibility for the claims and presentation.

\begin{thebibliography}{9}
\footnotesize
\raggedright

\bibitem{steinberger2012}
J.~P. Steinberger,
\emph{The lowest-degree polynomial with nonnegative coefficients divisible by
the $n$-th cyclotomic polynomial},
Electronic Journal of Combinatorics \textbf{19} (2012), no.~4, Paper P1.
\href{https://doi.org/10.37236/2755}{doi:10.37236/2755}.

\end{thebibliography}

\end{document}
